% ============================================================================
% 致谢
% 预估页数：1 页
% 写作要点：
%   - 标题"致 谢"两字之间空 2 格（西科大模板要求）
%   - 标题：黑体小一号，居中，段前段后 0.5 行
%   - 内容：宋体小四，固定行距 22 磅
%   - 致谢对象顺序通常为：导师 → 任课老师 → 同学 / 朋友 → 家人 → 学校
%   - 字数 400-600 字
% ============================================================================

\cleardoublepage
\pagestyle{thesisplain}
\chapter*{\centering 致\hspace{2em}谢}
\addcontentsline{toc}{chapter}{致谢}

\setlength{\parindent}{2em}
\songti\zihao{-4}

时光荏苒，四载光阴如白驹过隙。在即将完成本科毕业论文之际，谨向所有曾给予我关怀、教诲与帮助的人致以最诚挚的感谢。

首先，衷心感谢我的指导老师在毕业设计与论文写作全过程中给予的悉心指导。从选题论证、需求分析、系统设计，到功能实现、论文结构调整与格式规范，老师都提出了许多具体而中肯的意见，使我能够在较为复杂的工程实践中保持清晰的研究主线。特别是在人工智能技术与医学教育场景结合的过程中，老师反复提醒我注意系统目标、教学价值与工程可实现性之间的平衡，这对本文的完成具有重要帮助。

感谢西南科技大学计算机科学与技术学院各位老师四年来的教导与培养。专业课程学习为本课题的完成奠定了基础，软件工程、数据库系统、Web 开发、人工智能等课程中的知识在系统架构设计、数据建模、接口实现和智能问诊模块开发中都得到了实际应用。学院提供的学习环境与实践机会，也让我逐步形成了工程化分析问题和解决问题的能力。

感谢同学和朋友们在项目开发与论文撰写期间给予的帮助。系统测试、界面体验、论文表述和材料整理过程中，大家提出的建议帮助我发现了许多容易忽视的问题，使系统功能更加完整，论文表达也更加清晰。与同窗共同学习、讨论和调试的经历，是本科阶段非常宝贵的记忆。

感谢我的家人在求学期间给予的理解、支持与鼓励。正是他们在生活和精神上的陪伴，让我能够更加专注地完成学业和毕业设计。在遇到开发阻塞、论文修改和时间压力时，家人的支持使我能够保持耐心并坚持完成任务。

最后，感谢学校提供的学习平台和成长环境。毕业论文的完成既是本科阶段学习成果的一次集中检验，也是新的起点。今后我将继续保持严谨求实的学习态度，把在本课题中积累的工程实践能力和问题分析能力运用到后续学习与工作中。
